% NLP-BR-Sample.tex
% v1.1, released 10 Apr 2024
% Copyright 2024 Cambridge University Press

\documentclass[bookreview]{CUP-JNL-NLP}

\usepackage{amsmath}
\usepackage{graphicx}
\usepackage{natbib}
\ifpdf%
\usepackage{epstopdf}%
\else%
\fi
\usepackage{multirow}

\begin{document}
\label{firstpage}

\lefttitle{Book Review}
\righttitle{Book Review}

\jnlPage{\pageref{firstpage}}{\pageref{lastpage}}
\jnlDoiYr{2024}
\doival{10.1017/xxxxx}

\title{Book Review}

\maketitle

\bktext{Computational Linguistics, by Author. Cham, Springer Nature,
2024. ISBN 9123456789. IX + 22 pages.}

The layout book review design for the {\em Natural Language Processing} journal
has been implemented as a \LaTeX\ style file. The NLP style file is based
on the ARTICLE style as discussed in the \LaTeX\ manual. Commands which
differ from the standard \LaTeX\ interface, or which are provided in addition
to the standard interface, are explained in this guide. This guide is not a
substitute for the \LaTeX\ manual itself.

The \LaTeX\ document preparation system is a special version of the
\TeX\ typesetting program. \LaTeX\ adds to \TeX\ a collection of
commands which simplify typesetting by allowing the author to
concentrate on the logical structure of the document rather than
its visual layout.

\LaTeX\ provides a consistent and comprehensive document preparation
interface. There are simple-to-use commands for generating a table of
contents, lists of figures and/or tables, and indexes. \LaTeX\ can
automatically number list entries, equations, figures, tables, and
footnotes, as well as parts, chapters, sections and subsections.
Using this numbering system, bibliographic citations, page references
and cross references to any other numbered entity ({\it e.g.\ } chapter,
section, equation, figure, list entry) are quite straightforward.

The use of document class allows a simple change of style (or style option)
to transform the appearance of your document. The CUP NLP class file preserves
the standard \LaTeX\ interface such that any document which can be produced
using the standard \LaTeX\ ARTICLE style can also be produced with the
NLP style. However, the fonts (sizes) and measure of text is slightly different
from that for ARTICLE, therefore line breaks will change and it is possible
that equations may need re-setting.

\begin{contrib}
Author 1
Affiliation 1 Line 1,
Affiliation 1 Line 2
Email 1\\
Author 2
Affiliation 2 Line 1,
Affiliation 2 Line 2
Email 2
\end{contrib}

\begin{thebibliography}{}
  \bibitem[\protect\citeauthoryear{Akmajian and Lehrer}{Akmajian and Lehrer}{1976}]{akm76}
   \textbf{Akmajian and Lehrer A.} (1976). NP-like quantifiers and the
   problem of determining the head of an NP. {\it Linguistic
   Analysis\/} \textbf{2}, 295--313.
\end{thebibliography}

\label{lastpage}

\end{document}
